\vspace*{1.5cm}
\section{Resumen}

Se presenta un resumen del Trabajo Práctico 3, detallando los temas a tratar para ofrecer una lectura fluida, clara y precisa que permita comprender el desarrollo en su totalidad.

En primer lugar, se va a hablar sobre el problema propuesto, cómo fue el planteo de la resolución, cómo se logró llegar a que ésta misma pueda ser óptima para cualquier caso posible, las consecuencias temporales que genera eso y cómo empeorar la optimalidad nos puede traer una gran ventaja.

Después, se trata con profundidad el análisis de la complejidad de los problemas y cómo es posible categorizarlosrizarlos en conjuntos siempre y cuando cumplan condiciones polinómicas y/o exponenciales.

Se detalla la teoría necesaria sobre los algoritmos de \textbf{Backtracking}, algoritmos de \textbf{Programación Lineal} y algoritmos de \textbf{Aproximación}. Para ver estos aspectos aplicados, se incluye un breve ejemplo de la idea en general y su analogía a conceptos en el mundo real (toma de decisiones en la vida). Por otro lado, se van a detallar el diseño de los algoritmos implementados en \textit{pseudocódigo}, explicando anteriormente la lógica, cómo fue buscada y cómo fue pensada esa resolución. Adjuntándose al final la implementación realizada.

Se realiza un \textit{análisis} de cada algoritmo planteado, mediante el estudio de la complejidad y una demostración matemática para validar la optimalidad o cota del algoritmo, de forma clara y con ejemplos propuestos. Se incorporan múltiples experimentos empíricos de forma tal que se pueda visualizar gráficamente las medidas de tiempo sobre cierta cantidad de casos y poder comparar todos los algoritmos implementados.

Finalmente, se expone una conclusión sobre lo aprendido, la dificultad del problema, y la particularidad de estos algoritmos y su pensamiento.
